% s-180-p180-m003; original PDF 180, 181
\Aphorism{s-180-p180-m003}{270}
Wenn es ein quietistisches Wetter giebt, so haben wir’s heute. Gleichmässig grauer Himmel, ohne eine Spur Sonne, aber ebenso ohne Gewitterschwärze. Die Luft weder warm noch kühl, die Beleuchtung weder hell noch dunkel; beides neutral, auf dem Nullpunkt der Empfindung. Völlige Einsamkeit und Stille, die nicht als Stille fühlbar wird; ein fernes Kirchengeläut summt so eintönig dazwischen, dass ich mich erst besinnen muss, ob es noch tönt oder nicht. Dem matten Olivengrün fehlt heute jeder belebende Stich ins Silberne; der Berghang drüben sieht nach nichts aus, weder nach Nah noch Fern, nicht düster und nicht freundlich. Die Umwelt verzichtet, sich mir gegenüberzustellen, Relation zu bekennen, ein eigenes Gesicht zu machen; sie hebt sich gerade nur als Nicht-Ich von meinem Ich ab — — nichts weiter. Ich beginne etwas wie das reine eleatische Sein zu fühlen, dem jede sinnliche Färbung und Modification abgeht; es ist noch Leben, aber mit Abstraction von allem Lebendigen — nur eben noch unterschieden von der Unbeweglichkeit des absoluten Verharrens. Ich schlendere in den Bergwiesen umher, in einem Tempo mit dieser schleichenden Schnecke von Tag; jede rasche Action wäre hier eine tiefe Zwecklosigkeit, ein Zerreissen des milieu, ein durch Nichts herausgeforderter Widerspruch. Ich denke an nichts, höchstens an das Nichts; das Denkorgan ist zu schlaff, um sich zum Umspannen einer Begriffssphäre zu krümmen. Ich versinke in Ruhe, aber nicht als den Gegensatz des Bewegten; dieses Nirwana ist noch nie vom Sansara durchbrochen worden. Du Sein, das nur kein Nichtsein ist, du Ichts, vertreten durch ein nicht-Nichts, du Schlaf ohne Traum, du graues rinnendes Unding und Abstractum: wann bist du zu Ende, du — \Sperr{englischer Sonntag?}
\par

